\scanpage{6}\mspage{5}

À fortiori, \uline{le $A$-torseur}

\[
\underline{\mathrm{Alb}}^{\chi(\delta)}_{X/S}\overset A\times\underline{\mathrm{Pic}}^{\varphi(\delta)}_{B/S}\overset A\times(P^\delta\overset B\times(A,\widetilde{\varphi(\delta)}))
\]

\uline{est trivial} (en fait, avec une trivialisation canonique).

e) Pour exploiter ces considérations, il \deleted{\unsure{faudra,|faudrait}} encore, pour tout $\delta\in\mathrm{NS}(X/S)$, ``calculer'' $\varphi(\delta)$ (et $\chi(\delta)$ \ldots). Pour ceci, on \deleted{peut supposer} est ramené au cas $S=\mathrm{Spec}\,k$, \unsure{($k$|$k$|et $k$} un corps, alg. clos si on y tient\deleted{\unsure{)}}. Choisissons alors $p\in P^\delta(k)$, qui correspond à un faisceau inversible $(\mathcal L_g)_p=L$ sur $X$, dans la classe $\delta$.

\margin{$X\times P^\delta$}

\unsure{Le choix|La donnée} de $\deleted{\unsure{g}}p$ définit un isom

\[
P^\delta\simeq B
\]

et moyennant celui-ci, le faisceau $\mathcal L_g\mid X\timesP^\delta$ \deleted{$\unsure{X\times P^\delta}$} est transformé en \deleted{\unsure{$\mathcal L_g$ de $X\times$}} le faisceau $\mathcal L^{\unsure{0|\circ}}\deleted{\unsure{\otimes_X|\otimes}}\underline{\mathrm{pr}}_1^*(L_{\unsure{\delta|{}}})$, où $\mathcal L^\circ$ est le faisceau de Weil sur $X\times B$ [relatif au choix du point $g\in X(k)$) et où $L$ est \deleted{\unsure{de}}le faisceau inversible sur $X$ \deleted{défini} de classe $\delta$ déjà envisagé. Il s'agit donc de calculer

\[
\det{}^{*}(R\mathrm{pr}_{2*})(\mathcal L\otimes\mathrm{pr}_1^*(L))
\]
